\documentclass[11pt, a4paper]{article}
% Remove czech babel for english 
\usepackage[czech]{babel}
\usepackage[utf8]{inputenc}
\usepackage[T1]{fontenc}
\usepackage[left=2cm, top=3cm, text={17cm, 24cm}]{geometry}
\usepackage{times}
\usepackage[unicode, hidelinks, pdfpagelabels=false]{hyperref}
\usepackage{graphicx}
\usepackage{wrapfig,lipsum}

\graphicspath{ {./images/} }

% Custom properties
\newcommand\workType{Type of project}
\newcommand\class{Test class}
\newcommand\subject{\textbf{Theme or subject of theses}}
\newcommand\name{Title. Name of student}
\renewcommand{\refname}{Literatura}
% Change czplain to enplain for english citations
\newcommand\bibstylename{czplain}


\urlstyle{same}
\begin{document}

\begin{titlepage}
    \begin{center}
        \vspace*{\stretch{0.291}}
        \LARGE
        \class{} \\ 
        \Huge
        \subject{} \\
        \vspace{\stretch{0.091}}
        \LARGE
        \name{} \\
        \vspace{\stretch{0.518}}

    \end{center}

    \begin{wrapfigure}{R}{0.5\textwidth}
        \vspace{-33pt}
        \begin{center}
            \includegraphics[width=0.48\textwidth]{fai.eps}
        \end{center}
    \end{wrapfigure}

    {{\noindent \Large\workType{}} \\
    \\
    \Large\the\year{}}

\end{titlepage}

\tableofcontents

\newpage
\section{Book citation}

\uv{\emph{Typografie je skládání písmen a textu způsobem, který činí dokument čitelný, jasný a vizuálně přitažlivý}}.
Typografii můžeme chápat jako soubor pravidel, pro práci s~digitálním textem viz \cite{prakt:2000}. Právě dodržování těchto pravidel by nám mělo zajistit typografickou správnost výsledného textu.
Tyto pravidla ovšem nelze dodržovat bezhlavě. Každý text má svoje specifika a podle těch můžeme volit různá pravidla, kterých se budeme držet, mohli bychom tedy říci, že v~případě typografie vede k~správnému výsledku mnoho správných řešení.

Kromě dodržování typografických pravidel je vhodné při tvorbě dokumentu myslet i na obecné zásady tvorby dokumentu.

Nedílnou součástí typografie je písmo. Právě využití správných fontů a řezů písma může činit text čitelnější a daleko přitažlivější.

\newpage

\section{Web citation}
Písmo ačkoliv se to nemusí na první pohled zdát, má mnoho částí a vlastností viz \cite{typeface:2008}. Je samozřejmě také nedílnou součástí grafického 
designu a jsou to právě grafici, kteří nové fonty vytváří. Většina těchto návrhářů vytváří své fonty pro své zákazníky přímo a nebo 
vytváří fonty, které následně umístí na internetové obchody. Najdou se ovšem i takoví, kteří kromě zakázkové tvorby vytvářejí i písma, které dále publikují k~volnému užití, jako třeba Jonathan Hoefler.

Písma nejsou samozřejmě důležité pouze u~dokumentů, ale hrají obrovskou roli v~reklamě a také webových stránkách. V~těchto případech často slouží k~provokaci, která vyústí k~následnému kliknutí na odkaz, popřípadě jiný textový prvek.
Zajímavostí u~webových stránek je dělení písem na bezpečná a nebezpečná, kdy bezpečná písma jsou ta, u~kterých odpadá riziko, že je koncový uživatel uvidí jinak, než vývojář, který web vytváří.

\section{Article citation}
\lipsum[2-3]\cite{olsak:2014}

\newpage

\section{Master theses citation}
Pro vysvětlení toho, co je to \LaTeX musíme nejprve vědět, co je \TeX. {\TeX} je sázení systém, který narvhnul slavný informatik Donald Knuth v~roce 197. {\LaTeX} je také sázecí systém, stojí nad {\TeX}em a zjednodušuje s~ním interakci ,
vytvořil jej americký počítačový specialista Leslie Lamport na počátku 80. let. Od doby, kdy byly oba systémy vydány jsou hojně aktualizovány, aby uspokojovali potřeby dnešní doby.

Mnoho lidí bere {\LaTeX} jako velmi nepřátelský sázecí systém, který jim spíše práci stěžuje než urychluje, nicméně opak je pravdou, pokud se začátečník drží správného návodu viz \cite{sirucek:2007}, tak se s~{\LaTeX}em velice rychle sžije a práce se poté mnohonásobně zrychlí.

\newpage

\bibliographystyle{bibstyles/\bibstylename}
\addcontentsline{toc}{section}{\refname}
\bibliography{utb_template}
\end{document}